%%%%%%%%%%%%%%%%%%%%%%%%%%%%%%%%%%%%%%%%%%%%%%%%%%
\needspace{12\baselineskip}
\section{Java: Ausdrücke und Operatoren}
%%%%%%%%%%%%%%%%%%%%%%%%%%%%%%%%%%%%%%%%%%%%%%%%%%
\worum
Ein \textbf{Ausdruck} beschreibt eine Auswertungsvorschrift: Operatoren werden auf Operanden (Variablen, Literale, Methodenaufrufe) angewandt, und am Ende steht ein Wert. In welcher Reihenfolge die Operatoren ausgewertet werden, bestimmen ihre \textbf{Prioritäten}; mit runden Klammern lässt sich eine andere Reihenfolge erzwingen. Einige Operatoren haben \textbf{Nebeneffekte}: Sie liefern nicht nur einen Wert, sondern verändern nebenbei eine Variable. Genau das macht Ausdrücke wie \code{x = y++ + --z} zu einer beliebten Klausuraufgabe.

\begin{merke}[Die wichtigsten Operatoren]
\begin{itemize}[itemsep=0.1em]
\item \textbf{Arithmetisch:} \code{+ - * /} und \code{\%} (Rest der Division, Modulo).
\item \textbf{Inkrement und Dekrement:} \code{++} erhöht um 1, \code{--} verringert um 1.
	\begin{itemize}[itemsep=0em]
	\item Präfix \code{++x}: \emph{erst} ändern, \emph{dann} den neuen Wert verwenden.
	\item Postfix \code{x++}: \emph{erst} den alten Wert verwenden, \emph{dann} ändern.
	\end{itemize}
\item \textbf{Zuweisung:} \code{=} sowie die Kurzformen \code{+= -= *= /= \%=}. \ \code{a op= b} bedeutet \code{a = a op b}.
\item \textbf{Vergleich:} \code{== != < <= > >=}, das Ergebnis ist immer ein \code{boolean}.
\item \textbf{Logisch:} \code{\&\&} (und), \code{||} (oder), \code{!} (nicht), \code{\textasciicircum} (entweder oder). Bei \code{\&\&} und \code{||} wird der rechte Operand nur ausgewertet, wenn das Ergebnis noch nicht feststeht (\emph{kurze Auswertung}). Die Varianten \code{\&} und \code{|} werten immer beide Seiten aus.
\item \textbf{Bedingungsoperator:} \code{Bedingung ? Ausdruck1 : Ausdruck2} liefert \code{Ausdruck1}, falls die Bedingung wahr ist, sonst \code{Ausdruck2}.
\end{itemize}
\end{merke}

\begin{merke}[Prioritäten (von hoch nach niedrig)]
\begin{center}
\begin{tabular}{@{}cl@{\qquad\qquad}cl@{}}
1 & \code{++ -- !} Vorzeichen, Cast & 5 & \code{\&\&} \\
2 & \code{* / \%} & 6 & \code{||} \\
3 & \code{+ -} & 7 & \code{? :} \\
4 & \code{< <= > >=}, danach \code{== !=} & 8 & \code{= += -= *= /= \%=}
\end{tabular}
\end{center}
Operanden werden in Java immer \textbf{von links nach rechts} ausgewertet. Nebeneffekte eines linken Operanden sind also schon geschehen, wenn der rechte ausgewertet wird.
\end{merke}

\beispiele
\beispiel{Präfix und Postfix}
\begin{tcbraster}[raster columns=2, raster equal height, raster column skip=3mm]
\begin{javacode}[Ausschnitt]{Postfix}
int n = 5;
int m = n++;     // m = 5, n = 6
\end{javacode}
\begin{javacode}[Ausschnitt]{Präfix}
int n = 5;
int m = ++n;     // m = 6, n = 6
\end{javacode}
\end{tcbraster}

\beispiel{Einen Ausdruck Schritt für Schritt auswerten}
Welche Werte haben \code{x} und \code{y} nach \code{int x = 2; int y = x++ * 3 + --x;}?
\begin{center}
\begin{tabular}{@{}llc@{}}
\toprule
\textbf{Teilausdruck} & \textbf{Wert} & \textbf{x danach} \\
\midrule
\code{x++} & 2 (alter Wert) & 3 \\
\code{x++ * 3} & $2 \cdot 3 = 6$ & 3 \\
\code{--x} & 2 (neuer Wert) & 2 \\
\code{x++ * 3 + --x} & $6 + 2 = 8$ & 2 \\
\bottomrule
\end{tabular}
\end{center}
Ergebnis: \code{y = 8} und \code{x = 2}.

\beispiel{Zusammengesetzte Zuweisung mit Nebeneffekt}
Bei \code{a op= b} merkt sich Java \emph{zuerst} den alten Wert von \code{a}, wertet dann \code{b} aus und weist am Ende das Ergebnis zu.
\begin{javacode}[Ausschnitt]{Zuweisung}
int x = 4;
x += x++;        // 4 + 4 = 8.   Die Erhöhung auf 5 wird von der Zuweisung überschrieben.

int y = 3;
y *= ++y;        // 3 * 4 = 12.  Links zählt der alte Wert 3, rechts der neue Wert 4.
\end{javacode}

\beispiel{Bedingungsoperator}
\begin{javacode}[Ausschnitt]{Bedingungsoperator}
int max = (a > b) ? a : b;                 // das Maximum von a und b
int betrag = (a < 0) ? -a : a;             // der Betrag von a
String text = (n % 2 == 0) ? "gerade" : "ungerade";
\end{javacode}

\beispiel{Logische Operatoren und kurze Auswertung}
\begin{javacode}[Ausschnitt]{Logik}
int num1 = 2, num2 = 3;
System.out.println(num1 == 2 && num2 < 8);          // true
System.out.println(num1 != 2 || !(num2 < 2));       // true
System.out.println(num1 == 2 ^ num2 > 0);           // false (beide Seiten sind wahr)

int n = 0;
System.out.println(n != 0 && 10 / n > 2);           // false, es wird nicht durch 0 geteilt
\end{javacode}
Mit \code{\&} statt \code{\&\&} würde in der letzten Zeile auch die rechte Seite ausgewertet, und das Programm bräche mit einer \code{ArithmeticException} ab.

\aufgaben
\aufgabe{Inkrement und Dekrement}
Welche Werte haben die Variablen nach Ausführung der Ausschnitte? Überlegen Sie zuerst ohne Rechner.
\begin{tcbraster}[raster columns=3, raster equal height=rows, raster column skip=3mm, raster row skip=1mm]
\begin{javacode}[Ausschnitt]{a)}
int n = 7;
int m = n--;
\end{javacode}
\begin{javacode}[Ausschnitt]{b)}
int a = 4;
int b = --a * 3;
\end{javacode}
\begin{javacode}[Ausschnitt]{c)}
int x = 2;
int y = x++ + x++;
\end{javacode}
\begin{javacode}[Ausschnitt]{d)}
int i = 6;
int j = --i - i--;
\end{javacode}
\begin{javacode}[Ausschnitt]{e)}
int k = 5;
k -= k--;
\end{javacode}
\begin{javacode}[Ausschnitt]{f)}
int p = 3;
int q = ++p * p++;
\end{javacode}
\end{tcbraster}

\aufgabe[\K]{Programmverständnis I}
Gegeben seien die Variablenbelegungen \code{int x = 2; int y = 5; int z = 3;}

Welchen Wert haben die Variablen nach der Ausführung der einzelnen alternativen Zeilen? \textbf{Der Ausgangswert bei jeder Zeile sei die oben angegebene Belegung.}
\begin{center}
\renewcommand{\arraystretch}{1.5}
\begin{tabular}{|L{6.5cm}|C{1.5cm}|C{1.5cm}|C{1.5cm}|}
\hline
 & \textbf{x} & \textbf{y} & \textbf{z} \\ \hline
\code{z *= ++z;} & & & \\ \hline
\code{x = y-- - ++z;} & & & \\ \hline
\code{y -= (x < z) ? x : 2 * x;} & & & \\ \hline
\code{x = (y++ > 5) ? y : 3 * y;} & & & \\ \hline
\end{tabular}
\end{center}

\aufgabe[\K]{Programmverständnis II}
Gegeben seien die Variablenbelegungen \code{int x = 4; int y = 1; int z = 2;}

Welchen Wert haben die Variablen nach der Ausführung der einzelnen alternativen Zeilen? \textbf{Der Ausgangswert bei jeder Zeile sei die oben angegebene Belegung.}
\begin{center}
\renewcommand{\arraystretch}{1.5}
\begin{tabular}{|L{6.5cm}|C{1.5cm}|C{1.5cm}|C{1.5cm}|}
\hline
 & \textbf{x} & \textbf{y} & \textbf{z} \\ \hline
\code{z = y++;} & & & \\ \hline
\code{y = y++ + ++y;} & & & \\ \hline
\code{x += x *= z--;} & & & \\ \hline
\code{z -= (x-- - --y);} & & & \\ \hline
\end{tabular}
\end{center}

\aufgabe{Logische Ausdrücke}
Gegeben seien \code{int a = 5; int b = 10; boolean c = false;} Welchen Wert liefern die folgenden Ausdrücke?
\begin{multicols}{2}
\begin{enumerate}[label=\alph*)]
\item \code{a < b \&\& !c}
\item \code{a > b || b \% a == 0}
\item \code{a == 5 \textasciicircum\ b == 10}
\item \code{!(a < b) || c \&\& a > 0}
\end{enumerate}
\end{multicols}
\begin{enumerate}[label=\alph*), start=5]
\item Welche Werte haben \code{a} und \code{b} nach der Auswertung von \code{(a++ > 5) \&\& (b++ > 5)}? Was ändert sich, wenn \code{\&} statt \code{\&\&} verwendet wird?
\end{enumerate}

\aufgabe{Bedingungsoperator}
\begin{enumerate}[label=\alph*)]
\item Schreiben Sie die folgende Verzweigung als eine einzige Zuweisung mit dem Bedingungsoperator.
\begin{javacode}[Ausschnitt]{Verzweigung}
if (punkte >= 50) {
	ergebnis = "bestanden";
} else {
	ergebnis = "nicht bestanden";
}
\end{javacode}
\item Schreiben Sie ein Java-Programm, welches das Maximum dreier eingelesener Zahlen $a$, $b$ und $c$ ermittelt. Verwenden Sie zur Bestimmung den Bedingungsoperator.
\item Geben Sie einen Ausdruck an, der für eine Zahl \code{n} das Vorzeichen als $-1$, $0$ oder $1$ liefert.
\end{enumerate}

\aufgabe{Überlauf erkennen}
\begin{enumerate}[label=\alph*)]
\item Lesen Sie eine \code{int}-Variable \code{a} von der Konsole ein und weisen Sie sie mit einem Cast einer \code{byte}-Variablen \code{b} zu. Prüfen Sie mit dem Bedingungsoperator, ob der Wertebereich von \code{byte} ausgereicht hat, und geben Sie \code{b} sowie eine passende Meldung aus. Testen Sie mit 127, 128 und 300.
\item Lesen Sie zwei positive \code{int}-Zahlen ein und geben Sie ihre Summe aus. Prüfen Sie, ob es bei der Addition zu einem Überlauf gekommen ist. Testen Sie mit 2147483647 und 2.
\end{enumerate}

\aufgabe[\S]{Prioritäten}
\begin{enumerate}[label=\alph*)]
\item Setzen Sie Klammern so, dass die Auswertungsreihenfolge sichtbar wird, ohne das Ergebnis zu verändern: \\
\code{a + b * c > d || !e \&\& f} \qquad\qquad \code{x = y = z + 1} \qquad\qquad \code{a < b == c < d}
\item Welche Werte liefern \code{7 / 2 * 2.0}, \ \code{7 / 2.0 * 2}, \ \code{-7 / 2} und \code{-7 \% 3}? Begründen Sie.
\item Welchen Wert hat \code{i} nach der Anweisung \code{i = i++;}? Erklären Sie Schritt für Schritt.
\end{enumerate}

\ueben
\aufgabe{Zuweisungsoperatoren}
Welchen Wert hat \code{a} nach jeder Zeile? Die Zeilen werden nacheinander ausgeführt.
\begin{javacode}[Ausschnitt]{Nacheinander}
int a = 10;
a += 5;
a -= 3;
a *= 2;
a /= 5;
a %= 3;
\end{javacode}

\aufgabe{Noch mehr Inkrement}
\begin{tcbraster}[raster columns=3, raster equal height=rows, raster column skip=3mm, raster row skip=1mm]
\begin{javacode}[Ausschnitt]{a)}
int a = 1;
int b = a++ + a;
\end{javacode}
\begin{javacode}[Ausschnitt]{b)}
int a = 1;
int b = a + a++;
\end{javacode}
\begin{javacode}[Ausschnitt]{c)}
int x = 8;
int y = x-- / 2;
\end{javacode}
\begin{javacode}[Ausschnitt]{d)}
int x = 3;
x = x++ * 2;
\end{javacode}
\begin{javacode}[Ausschnitt]{e)}
int m = 5, n = 5;
boolean g = m++ == ++n;
\end{javacode}
\begin{javacode}[Ausschnitt]{f)}
int i = 2;
int j = i++ + i++ + i++;
\end{javacode}
\end{tcbraster}

\aufgabe{Bedingungsoperator auswerten}
Gegeben seien \code{int x = 7; int y = 3;} Welchen Wert liefern die Ausdrücke?
\begin{multicols}{2}
\begin{enumerate}[label=\alph*)]
\item \code{x > y ? x : y}
\item \code{x \% 2 == 0 ? x / 2 : 3 * x + 1}
\item \code{x < 5 ? "klein" : x < 10 ? "mittel" : "gross"}
\item \code{(x > y ? x : y) - (x < y ? x : y)}
\end{enumerate}
\end{multicols}

\aufgabe{Bedingungen formulieren}
Schreiben Sie die Aussagen als booleschen Ausdruck in Java.
\begin{enumerate}[label=\alph*)]
\item \code{x} liegt zwischen 1 und 10 (einschließlich).
\item Das Zeichen \code{c} ist ein Großbuchstabe.
\item \code{a} und \code{b} haben dasselbe Vorzeichen (beide positiv oder beide negativ).
\item Genau eine der beiden Zahlen \code{a} und \code{b} ist 0.
\item \code{n} ist durch 3 oder durch 5 teilbar, aber nicht durch beide.
\end{enumerate}

\aufgabe[\K]{Programmverständnis III}
Gegeben seien die Variablenbelegungen \code{int x = 1; int y = 3; int z = 2;}

Welchen Wert haben die Variablen nach der Ausführung der einzelnen alternativen Zeilen? \textbf{Der Ausgangswert bei jeder Zeile sei die oben angegebene Belegung.}
\begin{center}
\renewcommand{\arraystretch}{1.5}
\begin{tabular}{|L{6.5cm}|C{1.5cm}|C{1.5cm}|C{1.5cm}|}
\hline
 & \textbf{x} & \textbf{y} & \textbf{z} \\ \hline
\code{x += ++y;} & & & \\ \hline
\code{y = x++ * z--;} & & & \\ \hline
\code{z = (y > x) ? y - x : x - y;} & & & \\ \hline
\code{x = y-- - --y;} & & & \\ \hline
\end{tabular}
\end{center}
